\section{Newman evolution and the scalar generating field}\label{sec:newman}

\subsection{Entire source transform and Fourier transport}
Put $F_\lambda(u)=\ee^{\lambda u^2}\Phi(u)$. Evenness gives
\begin{equation}\label{eq:fullH}
 H_\lambda(z)=\frac12\int_\R F_\lambda(u)\ee^{\ii zu}\dd u,
 \qquad \widehat F_\lambda(k)=\sqrt{2/\pi}\,H_\lambda(k).
\end{equation}
For every fixed $a,b\ge0$, differentiation under this integral yields
\begin{equation}\label{eq:mixedH}
 \partial_\lambda^a\partial_z^bH_\lambda(z)
 =\frac12\int_\R u^{2a}(\ii u)^bF_\lambda(u)\ee^{\ii zu}\dd u.
\end{equation}
For $|\lambda|\le M$ and $|\operatorname{Im}z|\le B$, its absolute
value is at most $b_0J_{2a+b,2}(M,B)/2$. The same domination on complex
$\lambda$ compacts proves joint entire dependence. In particular,
\begin{equation}\label{eq:newmanPDE}
 \partial_\lambda^aH_\lambda=(-\partial_z^2)^aH_\lambda.
\end{equation}
This is an identity of the complete family, with no spectral truncation.

Finite Leibniz differentiation and Lemma~\ref{lem:tails} also imply
\begin{equation}\label{eq:Schwartzbound}
 |\partial_u^h(u^{2a}F_\lambda)(u)|
 \le C_{a,h,M}(1+|u|)^{2a+h}
 \ee^{Mu^2+|u|}x^{h+2}\ee^{-x}.
\end{equation}
For example one can generate derivatives of the Gaussian by
$R_0=1$, $R_{j+1}=\partial_uR_j+2\lambda uR_j$, so that
$\partial_u^j\ee^{\lambda u^2}=\ee^{\lambda u^2}R_j$;
the finite absolute coefficient sums determine $C_{a,h,M}$.
Thus $F_\lambda$, all polynomial insertions, and their finite parameter
derivatives are Schwartz, uniformly in each fixed seminorm on compact
real $\lambda$ intervals.

On this class the two operations are
\begin{equation}\label{eq:kineticvsnewman}
 \widehat{(-\partial_u^2)^af}(k)=k^{2a}\widehat f(k),\qquad
 \widehat{u^{2a}f}(k)=(-\partial_k^2)^a\widehat f(k).
\end{equation}
After an explicit spatial imprint $u=v$, a spinor encoding $f(v)\chi_n$
therefore evolves under free heat by the multiplier
$\ee^{-\tau(k^2+4n^2)}$. It does not evolve by
$-\partial_k^2$. The source law $\partial_\lambda F_\lambda=u^2F_\lambda$
is consequently different from the free heat or kinetic unitary law.

\subsection{The scalar bridge and its strip of convergence}
The radial theta geometry provides a second, scalar description
\cite{RadialCompanion,TNDParent}. Let $R=\ee^{2u}$ and define
\begin{equation}\label{eq:VG}
 V(R)=1+R^{-1}-\Theta(R),\qquad
 G_0(u)=\tfrac12\ee^u V(\ee^{2u}).
\end{equation}
For $u\ge0$, $G_0(u)=\ee^{-u}/2-p(u)$, and inversion makes $G_0$
smooth and even. In particular the apparent cusps in the corresponding
$|u|$ formula cancel. The potential $V$ is a scalar generating potential,
not an interaction added to $L_0$.

For Euclidean radius $R$, radial coordinate $u$ and angle $\omega$,
\begin{equation}\label{eq:measureconjugation}
 \begin{gathered}
 y=R\omega,\quad \dd y=R^2\dd R\dd\Omega,
 \quad \dd\mu=\frac{\dd u\dd\Omega}{4\pi}
              =\frac{\dd y}{8\pi R^3},\quad g_E=4R^2g_C,\\
 \Delta_y[R^{-1/2}f(u,\omega)]
 =-\tfrac14R^{-5/2}\Lsc f.
 \end{gathered}
\end{equation}
This is a differential conjugation with the displayed measure factors;
it is not a unitary equivalence between those Euclidean and cylinder
$L^2$ spaces. Inversion also gives
$V(R)=1-2R^{-1}Q(\pi/R^2)$, which extends smoothly through the Euclidean
origin. Distributionally the central singularities of $1/R$ and
$\Theta$ cancel in $V$; they cannot be dropped separately.

Only the angular-constant scalar channel is used by $G_\lambda$.
Write $L=1-\partial_u^2$. Its decaying inverse has kernel
$\ee^{-|u-w|}/2$, so
\begin{equation}\label{eq:Glambda}
 LG_\lambda=8F_\lambda,\qquad
 G_\lambda=8L^{-1}F_\lambda
 =4\int_\R\ee^{-|u-w|}F_\lambda(w)\dd w.
\end{equation}
This is the unique solution decaying at both ends; a homogeneous
combination of $\ee^u,\ee^{-u}$ cannot decay at both ends unless it is zero.
It is positive, even, smooth and Schwartz. Splitting the integral at $w=u$
gives the more precise asymptotic
\begin{equation}\label{eq:Gtail}
 G_\lambda(u)\sim c_\lambda\ee^{-|u|},\qquad
 c_\lambda=4\int_\R\ee^wF_\lambda(w)\dd w=8H_\lambda(\ii)>0.
\end{equation}

Integration by parts in \eqref{eq:Glambda} is valid initially only in
$|\operatorname{Im}z|<1$ and yields
\begin{equation}\label{eq:Gobs}
 16H_\lambda(z)=(1+z^2)\int_\R G_\lambda(u)\ee^{\ii zu}\dd u.
\end{equation}
To justify all boundary terms, convolution gives for every fixed
source/spatial derivative and $0\le a<1$ the weighted supremum estimate
\[
 \norm{\ee^{a|u|}\partial_\lambda^p\partial_u^qG_\lambda}_\infty
 \le\frac{8}{1-a}
 \norm{\ee^{a|u|}\partial_u^q(u^{2p}F_\lambda)}_\infty.
\]
On a compact substrip choose $|\operatorname{Im}z|<a<1$.
The remaining exponential decay pays the integrals and endpoint limits.
By \eqref{eq:Gtail}, the raw integral in \eqref{eq:Gobs} is not absolutely
convergent on the strip boundary. At $z=\pm\ii$ it is invalid to multiply
that divergent integral by zero. The entire continuation is instead
defined by \eqref{eq:fullH}. At $\lambda=0$, $c_0=1/2$ and
$H_0(\pm\ii)=1/16$, which check both normalizations.

\subsection{The transported generator and heat comparison}
On the invariant Schwartz class define
\begin{equation}\label{eq:transport}
 \Bcal=L^{-1}X^2L,\qquad
 \Bcal f=u^2f+L^{-1}(4uf'+2f).
\end{equation}
The second identity follows from
$L(u^2f)=u^2Lf-4uf'-2f$. Finite powers telescope on this class:
\begin{align}
 \partial_\lambda^pG_\lambda
 &=\Bcal^pG_\lambda=L^{-1}X^{2p}LG_\lambda,\label{eq:Bpower}\\
 \widehat{\Bcal^pf}(k)
 &=\frac{(-\partial_k^2)^p[(1+k^2)\widehat f(k)]}{1+k^2}.
\end{align}
The readout \eqref{eq:Gobs} then gives
$(-\partial_z^2)^pH_\lambda$. Its derivatives act on the factor
$1+z^2$ as well as on the integral. Replacing $\Bcal$ by $u^2$ acting
directly on $G$ would omit the product-rule terms. Indeed
$\ee^{\delta u^2}G_\lambda$ is not integrable for $\delta>0$,
whereas $\ee^{\delta u^2}F_\lambda$ is.

Let $\mathcal T$ denote the normalized observation in \eqref{eq:Gobs}.
The distinction between the two scalar heat generators is
\begin{equation}\label{eq:heatcompare}
 \begin{aligned}
 \mathcal T(\ee^{\tau\Delta_C}G_\lambda)(z)
   &=\ee^{-\tau z^2}H_\lambda(z),\\
 \mathcal T(\ee^{-\tau\Lsc}G_\lambda)(z)
   &=\ee^{-\tau(1+z^2)}H_\lambda(z).
 \end{aligned}
\end{equation}
These hold first in the paid strip and then by continuation of the
response; the second carries the extra $\ee^{-\tau}$
\cite[Section 5]{TNDSharedReview}. They do not describe the spinor
auxiliary response, which has its own angular factor in
\eqref{eq:auxintro}.

Finally, multiplication by $\ee^{\delta X^2}$ is unbounded on
unrestricted $L^2(\R)$ for $\delta>0$: normalized functions supported
in $[R,R+1]$ have image norm at least $\ee^{\delta R^2}$.
The canonical family is nevertheless defined for all finite real
$\lambda$ by Lemma~\ref{lem:tails}. This excludes a bounded realization
by that particular unrestricted multiplier, not state-specific
preparation, conditioning, or an enlarged model.

\section{Normalized sources and signed observations}\label{sec:normalized}

For real $\lambda$, use the full-line normalization
\begin{equation}\label{eq:fstate}
 Z_\lambda=\int_\R F_\lambda(v)\dd v=2H_\lambda(0),\qquad
 f_\lambda(v)=Z_\lambda^{-1/2}\ee^{\lambda v^2/2}\sqrt{\Phi(v)}.
\end{equation}
The use of $v$ supplies the spatial imprint. On $|\lambda|\le M$,
$0<Z_{-M}\le Z_\lambda\le Z_M<\infty$.

\begin{lemma}[Square-root regularity and complete angular preparation]
\label{lem:state}
Every fixed Schwartz seminorm of $f_\lambda$ and every fixed mixed
parameter/spatial derivative is bounded on a compact real $\lambda$
interval. Define the complete angular state
\begin{equation}\label{eq:R}
 R=\frac{W\ee^{-\pi A}}{Q(\pi)}
  =\sum_{n\ge1}\frac{\ee^{-\pi n^2}}{4nQ(\pi)}P_n,
 \quad \Tr R=1,
 \quad \mu_b:=\Tr(A^bR)=\frac{\sum_n n^{2b}\ee^{-\pi n^2}}{Q(\pi)}<\infty.
\end{equation}
For $\rho^I_\lambda=|f_\lambda\rangle\langle f_\lambda|\otimes R$,
the quadratic-form sandwich $L_0^p\rho^I_\lambda L_0^p$ is positive
trace class for each fixed $p\ge0$.
\end{lemma}
\begin{proof}
Schwartz regularity of a positive function alone does not prove it for
its square root. Here the matching lower and upper bounds in
\eqref{eq:twosidedbounds} give
$|\Phi^{(j)}/\Phi|\le(b_j/\kappa)x^j$.
Put $a_j=\Phi^{(j)}/\Phi$, so
$a_j'=a_{j+1}-a_ja_1$. If $h=\sqrt\Phi$, induction gives
$h^{(j)}=hS_j$, where
\begin{equation}\label{eq:sqrtRec}
 S_0=1,\qquad
 S_{j+1}=\tfrac12a_1S_j+
 \sum_{\ell=1}^j\frac{\partial S_j}{\partial a_\ell}
 (a_{\ell+1}-a_\ell a_1).
\end{equation}
Every monomial has weighted degree $j$ when $a_\ell$ has weight $\ell$.
Consequently, for finite $c_j$,
\[
 |h^{(j)}(v)|\le c_j\ee^{|v|/2}x^{j+1}\ee^{-x/2}.
\]
Multiplication by $\ee^{\lambda v^2/2}$ and finite Leibniz expansion
preserve all fixed Schwartz bounds. Derivatives of $Z_\lambda^{-1/2}$
are bounded using its positive lower bound and the paid integrals
$Z_\lambda^{(a)}=\int v^{2a}F_\lambda\dd v$. Smooth evenness and strict
positivity handle the origin.

For the complete angular state, Tonelli and the finite binomial theorem
give the exact trace payment
\begin{equation}\label{eq:tracemoments}
 \Tr(L_0^p\rho^I_\lambda L_0^p)
 =\sum_{j=0}^{2p}\binom{2p}{j}4^{2p-j}\mu_{2p-j}
 \norm{P^jf_\lambda}_2^2<\infty,\qquad P=-\ii\partial_v.
\end{equation}
The interpretation is the convergent positive sum of the weighted
ensemble rank-one vectors. It does not put every vector in the density's
closed support into $\Dom L_0^p$.
\end{proof}

All angular tails in \eqref{eq:tracemoments} are Gaussian series tails.
The radial tails, for any fixed nonnegative $j,N$ and $K>0$, satisfy
\[
 \int_{|k|>K}|k|^{2j}|\widehat f_\lambda(k)|^2\dd k
 \le K^{-2N}\norm{P^{j+N}f_\lambda}_2^2.
\]
Thus the full spectral moments, rather than only finite-mode states,
are controlled \cite[Section 3]{TNDDirect}.

The normalized Newman tangent and signed observations are
\begin{align}
 \partial_\lambda f_\lambda
  &=\tfrac12(X^2-\langle X^2\rangle_\lambda)f_\lambda,\label{eq:tangent}\\
 \frac{H_\lambda(k)}{H_\lambda(0)}
  &=\langle\cos(kX)\rangle_\lambda,\qquad
 \frac{H'_\lambda(k)}{H_\lambda(0)}
  =-\langle X\sin(kX)\rangle_\lambda.\label{eq:collisionobs}
\end{align}
Expectations use the probability density $|f_\lambda|^2=F_\lambda/Z_\lambda$.
The second observable is unbounded, but its expectation is paid by the
finite moments. A real multiple zero is simultaneous vanishing of the
two signed expressions in \eqref{eq:collisionobs}; a nondegenerate double
zero also has $H''_\lambda(k)\ne0$.
This is cancellation in observations, not coincidence of angular energy
levels. The Fourier transform of the density, the Fourier amplitude of
$f_\lambda$, and its momentum probability density are different objects.
